\chapter{Non-linear PDE}
\section{Selected Examples}


\section{The Fixed Point Theorem}
\begin{theorem}
Given $X$ a Banach space and $M \subset X$ \textbf{closed} and \textbf{convex}. If $T : M \rightarrow M$ is continuous, then $T$ admits at least one fixed point. 	
\end{theorem}


\subsection{Calculus of Variations}
\begin{definition}[Null Lagrangian]
	The Lagrangian is called \textbf{null-Lagrangian} if the corresponding \textbf{Euler-Lagrange function} is satisfied by all $C^2$-functions.
\end{definition}


\subsection{The Brouwer fixed point theorem}
The first tool in the toolbox of fixed point theorem is the Brouwer fixed point theorem. It states that in the finite dimensional space, continuous map on a closed ball admits at least 1 fixed point. Although we aren't able to apply it directly in most of cases since the functional space is usually infinite dimensional, it provides a off-hand tool in constructing finite dimensional existence. 

\begin{theorem}[Brouwer fixed point theorem(finite dimensional)]
	Given a continuous function defined on a closed ball in the finite dimensional real space, the function admits \textbf{at least} one fixed point. \par 
	Given $T: \bar{B_R(0)} \rightarrow \bar{B_R(0)}$ continuous with $B_R(0) \subset \mathbb{R}^m$, $T$ has at least one fixed point. 
\end{theorem}
\begin{remark}
The result is wrong in the infinite dimensional case. 	
\end{remark}

\begin{lemma}[no-retraction principle]
	There exists no continuous map $ w: B \rightarrow \partial B$, such that $w(x) = x$ for all $x \in \partial B$. 
\end{lemma}



\subsection{The Schrauder Fixed Point Theorem}
When we prove existence with fixed point theorem in the infinite dimensional case, Brouwer is no longer applicable. 

\subsubsection{Nonlinear compact operator }
\begin{proposition}[Approximation of compact operator by finite dimensional operator]
Given $M$ bounded $T: M \rightarrow Y$, then the followings are equivalent
\begin{enumerate}
	\item T is compact
	\item $\forall n \in \mathbb{N} \exists T_n : M \rightarrow Y: T_n(M) \subset Y_n$ with $\dim (Y_n) < \infty$, $T_n \overset{uniform}{\rightarrow} T$. 
\end{enumerate}
\end{proposition}
\begin{proof}\quad\\
\textbf{2)$\Rightarrow$ 1)}: \\
\textbf{i).} $T$ is compact. By uniform convergence of the operator we have
\begin{align*}
	\forall \epsilon > 0: \sup_{x \in M} \|Tx - T_{n_0}x \| \le \frac{\epsilon}{3} \quad \texttt{for $n_0$ sufficiently large}
\end{align*}
By the compactness of $T_{n_0}$ we have
\begin{align*}
	\forall \epsilon  \exists N \in \mathbb{N}: T_{n_0}(M) \subset \bigcup_{i = 1}^N B_{\frac{\epsilon}{3}} (x_i)
\end{align*}
therefore $\forall y \in M$
\begin{align*}
	\|T( y) - T(x_i)\|_Y \le \|T(y) - T_{n_0} (y)\|_Y + \|T_{n_0}(y) - T_{n_0}(x_i)\|_Y + \|T_{n_0}(x_i) - T(x_i)\|_Y
\end{align*}
on the RHS the first and the third are $\le \frac{\epsilon}{3}$ by uniform convergence and the second by the compactness of $T_{n_0}$. \\
\textbf{2)$ \Rightarrow$ 1)} 
\end{proof}



\begin{theorem}[Schrauder Version I]
	Given $X$ a Banach space, $M \subset X$ convex and compact, $T: M \rightarrow M$ continuous. Then $T$ admits a fixed point. 
\end{theorem}

\begin{lemma}[Mazur's lemma]
	Convex hull of precompact set is compact. 
\end{lemma}

\begin{theorem}[Schrauder Version II]
	Given $X$ a Banach space, $M \subset X$ convex, closed and bounded. The map $T: M \rightarrow M$ is compact. Then $T$ admits a fixed point. 
\end{theorem}

\subsection{An application to ODEs}
\begin{theorem}[Peano]
Given $(t_0, y_0) \in \mathbb{R} \times \mathbb{R}^n$, $Q:= \{(t, y): |t - t_0| \le a, |y - y_0| \le b\}$, let $f: Q \rightarrow \mathbb{R}^n$ be continuous and uniformly bounded, i.e., $|f(t, y)| \le K$ for some $K \in \mathbb{R}$. Then the integral equation
\begin{align*}
	y(t):= y_0 + \int_{t_0}^t f(s, y(s)) ds 
\end{align*}
admit a local solution.	
\end{theorem}

\begin{corollary}[Peano existence theorem in ODE]
	
\end{corollary}


\subsection{Application of Schrauder}



\section{Nonlinear Poisson Problems}
\begin{lemma}[Estimate of the inverse Dirichlet Laplacian]
Consider the bounded linear map on $\Omega$ with $\partial\Omega \in C^{1, 1}$
\begin{align*}
	\Delta_D: H^1_0(\Omega) \cap H^2(\Omega) &\rightarrow L^2\\
	u &\mapsto -\Delta u
\end{align*}
The operator is an isomorphism therefore we may define the inverse
\begin{align*}
	A_D: L^2 &\rightarrow H^1_0 \cap H^2\\
	g &\mapsto (-\Delta_D)^{-1}(g)
\end{align*}
and we have the estimate
\begin{align*}
	\|A_D\|_{H^2} \le \|g\|_{L^2}\\
	\|A_D\|_{H^1_0} \le \|g\|_{L^2}
\end{align*}
\end{lemma}


\subsection{Solvability with Schauder}
The PDE under consideration is the \textbf{nonlinear Poisson equation}.       
\begin{align*} 
\textbf{(NP)}\qquad
		- \Delta u = f(u) &\texttt{ on } \Omega\\
		u  = 0 &\texttt{ on } \partial \Omega 
\end{align*}           
\begin{theorem}[Well-posedness of the nonlinear Poisson Problem]
	Observe the non-linear elliptic system:
	\begin{align*}
		- \Delta u = f(u) &\texttt{ on } \Omega\\
		u  = 0 &\texttt{ on } \partial \Omega 
	\end{align*}
	If $f \in L^\infty(\Omega) \cap C^0(\Omega)$, then the system admits a weak solution $u \in H^1(\Omega)$. Moreover, if $\partial\Omega \in C^{1, 1}$, $u \in H^2(\Omega)$.
\end{theorem}
\begin{proof}
Define the map:
\begin{align*}
	T: L^2(\Omega) \rightarrow L^2(\Omega), u \mapsto A_D(f(u))
\end{align*}
we would like to apply the Schrauder to see that it admits a fixed point. 
	\quad\\
	\textbf{Step1:  $T$ is well-defined}: By the continuity of $f$ we know
	\begin{align*}
		\forall u \in L^2: f(u) \in L^\infty \in L^2
	\end{align*}
	Moreover we know the inverse Laplace operator satisfies
	\begin{align*}
		 \|A_D(f(u))\|_{L^2} \le \|A_D(f(u))\|_{H^1_0} \le \|f(u)\|_{L^2} < \infty
	\end{align*}
	since $f \in L^\infty$. Therefore $T$ is well-defined.\\
	\textbf{Step2: $\exists R > 0: T(B_R(0)) \subset \overline{B_R(0)}$, moreover, $B_R(0)$ is precaompact in $H^1_0(\Omega)$. } Observe $\forall u \in L^2(\Omega):$
	\begin{align*}
		\|T(f(u))\|_{L^2} \le \|A_D(\|f\|_\infty)\|_{L^2} 
		\le \|f\|_\infty \|A_D(1)\|_{L^2} \le \|f\|_\infty C_p |\Omega|^{\frac{1}{2}}: = R
	\end{align*}
	Therefore we have 
	\begin{align*}
		T(B_R(0)) \subset B_R(0)
	\end{align*}
	By Rellich we know $B_R(0)$ relative compact in $H^1_0(\Omega)$. 
	therefore we know $T(B_R(0))$ is precompact.
	\\
	\textbf{Step3: $T$ is continuous.} Given $u_n \rightarrow u$ in $L^2$, we know up to a subsequence 
	\begin{align*}
		u_k  \overset{a.s.}{\rightarrow} u
	\end{align*}
	observe
	\begin{align*}
		\|T(u_k) - T(u)\|_{L^2} 
		\overset{\texttt{linearity of $A_D$}}{=} \|A_D(f(u) - f(u_k))\|_{L^2} 
		\overset{lemma}{\le} \|f(u_k) - f(u)\|_{L^2} 
		\overset{\texttt{DCT and $f \in C^0$}}{\rightarrow} 0
	\end{align*}
	\\
	In conclusion, $T$ is compact on a convex and compact ball $B_R(0)$ in $H^1_0$, therefore $T$ admits a fixed point and the elliptic system admits a solution.
\end{proof}


\subsection{A priori estimate with weak maximum principle}
\begin{definition}
	for $u \in W^{1, p}(\Omega)$, we say $u \ge 0$ on $\partial\Omega$ if $u_ \in W^{1, p}_0(\Omega)$.
\end{definition}

\begin{proposition}
	Let $u \in H^1(\Omega)$ and assume
	\begin{enumerate}
		\item $u \ge 0$ on $\partial\Omega$
		\item $-\Delta u \ge 0$ in the weak sense, i.e., 
		\begin{align*}
			\int_\Omega \nabla u \cdot \nabla v dx \ge 0 \qquad \forall v \in H^1_0(\Omega)
		\end{align*}
		Then we have $u \ge 0$ a.e.
	\end{enumerate}
\end{proposition}

\begin{example}[Exponential diffusion]
Observe the system
\begin{align*}
	-\Delta u = e^{-u} &\texttt{ in } \Omega\\
	u = 0 &\texttt{ on } \partial \Omega
\end{align*}
Notice we have on $\partial\Omega$ $u \ge 0$ and the weak form 
\begin{align*}
	 \int_\Omega \nabla u \nabla \varphi = - \int_\Omega \Delta u \varphi  \ge 0 \quad \forall \varphi \in H^1_0(\Omega) \ge 0 
\end{align*}
By the weak maximum principle we have
\begin{align*}
	u \ge 0 \texttt{ a.e. in  } \Omega 
\end{align*}
we may truncate the diffusion term $f$ by
\begin{align*}
	\tilde{f} &= f \quad u \ge 0\\
	&= 1 \quad u \le 0 
\end{align*}
we have $f \in L^\infty \cap C^0(\mathbb{R})$. By the theorem in the last chapter we get a solution
\begin{align*}
	u \in H^1_0(\Omega)
\end{align*}
Moreover, the solution is unique since $f$ is non-increasing. Assuming $u_1, u_2 \in H^1_0(\Omega)$ are 2 solution of the diffusion equation, by linearity $u_1 - u_2$ is also a solution, we have
\begin{align*}
	\int_\Omega |\nabla(u_1 - u_2)|^2 dx = \int_\Omega (f(u_1) - f(u_2))(u_1 - u_2) \le 0 
\end{align*}
\end{example}





\section{Variational Inequalities \& Monotone Operators}
Differentiation in Banach space. 
\subsection{Variational Inequality}
\textbf{Motivation}:
\begin{align*}
	M:= \{v \in H^1_0(0, 1), v(0) = a, v(1) = b \}
\end{align*}
define the functional 
\begin{align*}
	I(v) = \int_I \sqrt{1 + v'} dx 
\end{align*}
want to find
\begin{align*}
	\min_{ v \in M} I(v) 
\end{align*}
Let $\varphi \in H^1_0(0, 1)$, we know $u$ is a minimiser of $I(v)$ if 
\begin{align*}
	\frac{d}{d\epsilon} I(u + \epsilon \varphi)|_{ \epsilon = 0}   = \int_I \bigg(\sqrt{1 + |u'|^2} \bigg)^{-1} u' \varphi' dx := \langle F(u), \varphi \rangle  
\end{align*}

\begin{remark}
If $u \in H^2(0, 1)$, then we have
\begin{align*}
	-\int_I (\frac{u'}{\sqrt{1 + |u'|^2}})' \varphi = 0 \quad \varphi \varphi \in H_0^1(\Omega)
\end{align*}
therefore we know $u'$ is constant and the minimal curve is a straight line. 	
\end{remark}
The \textbf{Ellipticity} is clear in this case.  

\quad\\
\textbf{The obstacle problem:} Let $h \in C^1(I, \mathbb{R})$ we define
\begin{align*}
	M:= \{v \in H^1_0(0, 1), v(0) = a, v(1) = b, v \ge h \}
\end{align*}
The Gateaux derivative cannot be expected anymore, on the other hand we observe the variational inequality:
\begin{align*}
	\langle F(u), u - v \rangle  \le 0 \qquad \forall v \in M 
\end{align*}
on a closed and convex set $M$. 



\quad\\
\textbf{General setting}
\begin{enumerate}
	\item $X$ is reflexive, separable Banach space with the dual space $X^*$. 
	\item $M \subset X$ closed, non-empty, convex. 
\end{enumerate}

consider the variational problem:
\begin{align*}
	I(v) = \int_a^b L( \nabla u, u, x) dt 
\end{align*}

The most simple example would be the shortest curve between 2 points
\begin{align*}
	I(v) := \int_a^b \sqrt{1 + u'(x)} dx  
\end{align*}
and consider the minimisation problem:
\begin{align*}
	\min_{v \in M} I(v) \qquad M := \{ v \in H^1([0, 1]) | v(0) = a, v(1) = b \}
\end{align*}
$M$ is well-defined by the classical embedding result.

\quad\\
\textbf{Assumptions on $F$:}\\
\textbf{F1:} $F$ maps bounded set of in $X$ to bounded set in $X^*$.\\
\textbf{F2:} Coercivity: There exists $u_0 \in M: \langle F(u), u - u_0 \rangle / \|u - u_0\| \rightarrow \infty$ as $\|u \| \rightarrow \infty$. \\
\textbf{F3:} Given $u_n \rightharpoonup u, F(u_n) \rightharpoonup u^*$, we have
\begin{align*}
	\langle u^*, u \rangle \le \liminf_{n \rightarrow } \langle F(u_n), u_n \rangle 
\end{align*}
Furthermore, if $\langle u^*, u \rangle = \liminf \langle F(u_n), u_n \rangle$, it holds
\begin{align*}
	\langle F(u) - u^*, u - v \rangle \le 0  \quad \forall v \in M
\end{align*}
\begin{remark}
\quad\\
\textbf{1).} The coercivity condition endowed $u$ with boundary growth condition.\\
\textbf{2).} \textbf{F3}) gives an upper bound for the variational inequality. \\
\textbf{3).} For \textbf{F2} it is essential to have $u_0 \in M$.
\end{remark}

\begin{example}[Counterexample for the coercivity condition]
\end{example}



\begin{theorem}
	If the operator $F$ satisfies \textbf{(F1) - (F3)}, then $F$ satisfies the variational inequality
	\begin{align*}
		\langle F(u), u - v \rangle \le 0
	\end{align*} 
\end{theorem}
\begin{proof}
	By Galerkin method. 
\end{proof}


\subsection{the Monotone Operators}
\begin{definition}[The monotone operator]
	An operator $A: M \rightarrow X^*$ is a \textbf{monotone operator} if 
	\begin{enumerate}
		\item $\langle A(x) - A(y), x - y \rangle \ge 0 \quad \forall x, y \in M$
		\item \textbf{hemi-continuous:} $\forall w \in X \forall v, u \in M$:
		\begin{align*}
			t \mapsto \langle A(u + t(v - u)), w \rangle 
		\end{align*}
		is continuous on $[0, 1]$.
		\item 
	Moreover, the operator is \textbf{strictly monotone} if 
	\begin{align*}
		\langle A(u) - A(v), u - v \rangle > 0 \qquad \forall u, v \in M
	\end{align*}
	\textbf{strongly monotone} if 
	\begin{align*}
		\exists c_0 \forall u, v \in M : \langle A(u) - A(v), u - v \rangle \ge c_0 \|u - v \|^2 
	\end{align*}
	\end{enumerate}
\end{definition}
hemi-continuity can be seen as one-sided continuity. It ensure the cone property of the monotone operator. As we have seen in the case of the operator $F$, by taking the convex combination of $u$ and $v$ the non-positivity of the operator can be extended from the cone corner $u$ to the direction of $v$. The hemi-continuity has the same function to preserve the cone property w.r.t. the nonlinearity $A$. 

\begin{lemma}[Minty]
	Given a monotone operator $F$, for $u \in M$ and $u^* \in X^*$, the following conditions are equivalent
	\begin{enumerate}
		\item $ \langle F(u) - u^*, u - v \rangle \le 0 \quad \forall v \in M$
		\item $\langle F(v) - u^*, u - v \rangle \le 0 \quad \forall v \in  M$.  
	\end{enumerate}
\end{lemma}
\begin{proof}
	\quad\\
	$\Rightarrow$: 
	\begin{align*}
		\langle F(v) - u^*, u - v \rangle = 
		\langle F(u) - u^*, u - v \rangle - \langle F(u) - F(v), u - v \rangle \le 0
	\end{align*}
	$\Leftarrow$: Define $v_\epsilon:= u + \epsilon( v - u)$ for $\epsilon \in [0, 1]$,  we know by assumption
	\begin{align*}
	 	0 \ge \langle F(v_\epsilon) - u^*, u - v_\epsilon \rangle = \langle F(v_\epsilon) - u^*, \epsilon (u - v) \rangle 
	 	\Rightarrow \langle F(v_\epsilon) -  u^*, u - v\rangle \le 0
	\end{align*}
	Letting $\epsilon \rightarrow 0$ we have by the hemi-continuity:
	\begin{align*}
		\langle F(u) - u^*, u - v \rangle \le 0
	\end{align*}
\end{proof}

\begin{lemma}
	Given a monotone operator $A$ we have
	\begin{enumerate}
		\item $A$ satisfies the lower semi-continuous condition \textbf{(F3)}.  
		\item If the operator is strictly monotone, then the variational inequality
				\begin{align*}
			\langle F(u), u - v \rangle \le 0 \quad \forall v \in M
		\end{align*}
		exists at most one solution. 
		\item If the operator is strongly monotone, then the coercivity condition holds w.r.t. all $u_0 \in C$
	\end{enumerate}
\end{lemma}
\begin{proof}
\quad\\
\textbf{1)}: Observe
\begin{align*}
	&0 \le \langle F(u) - F(u_n), u - u_n \rangle = \langle F(u), u - u_n \rangle + \langle F(u_n), u_n \rangle - \langle F(u_n), u \rangle \\
	&\Rightarrow \langle F(u_n) , u \rangle \le \langle F(u_n), u_n \rangle \\
	&\Rightarrow \langle u^*, u \rangle \le \liminf_{n \rightarrow \infty} \langle F(u_n), u_n \rangle \rightharpoonup \langle u^*, u_n \rangle \rightharpoonup \langle u^*, u \rangle 
\end{align*}
therefore
\begin{align*}
	\langle u^*, u \rangle = \liminf \langle F(u_n), u_n \rangle 
\end{align*}
Furthermore, by monotonicity
\begin{align*}
	&0 \le \langle F(u_n) - F(v), u_n - v \rangle = \langle F(u_n), u_n \rangle - \langle F(u_n), v \rangle  - \langle F(v), u_n - v \rangle\\
	&\Rightarrow
	\langle F(v), u_n - v \rangle - \langle F(u_n), u_n - v\rangle \le \langle F(u_n), u_n \rangle - \langle F(u_n), u_n \rangle\\
	&\Rightarrow 
	\langle F(v) - u^*,  u - v \rangle \le \liminf \langle F(u_n), u_n \rangle - \langle u^*, u \rangle  = 0
\end{align*}
By the Minty's trick we have
\begin{align*}
	\langle F(u) - u^*, u - v \rangle \le 0 \qquad \forall v \in M
\end{align*}
\textbf{2):} Suppose $u_1, u_2$ satisfies $\langle F(u), u - v \rangle \le 0 \quad \forall v \in M$, then  we have
\begin{align*}
	\langle F(u_1) - F(u_2), u_1 - u_2 \rangle \le 0
\end{align*}
which contradicts the strict monotonicity. \\
\textbf{3)}
\begin{align*}
	\langle F(u_0), u_0 - u \rangle \slash \|u_0 - u\|
	&= \langle F(u_0) - F(u), u_0 - u \rangle  \slash \|u - u_0 \| + \langle F(u), u_0 - u \rangle \slash \|u - u_0\|\\
	&\ge c_0 \|u - u_0\| + C
\end{align*}
where $c_0$ is the strong monotone constant of $F$ and $C$ depends on the operator norm of $F(u)$. We can conclude that
\begin{align*}
	\frac{\langle F(u_0), u_0 - u \rangle }{\|u_0 - u\|} \overset{\|u_0 -u\| \rightarrow \infty}{\longrightarrow} \infty
\end{align*}
\end{proof}


\subsection{Proof of the existence theorem}
\begin{proposition}
Given $X$ a Hilbert space, if $\dim X < \infty$, $F$ satisfies \textbf{(F1)} - \textbf{(F3)}, then there exists $u \in M$ with
	\begin{align*}
		\langle F(u), u - v \rangle \le 0 \quad \forall v \in M
	\end{align*} 
\end{proposition}
\begin{proof}
\quad\\
\textbf{Step 1: In a neighbourhood $B_\epsilon M$ of $M$:}\\
\textbf{Step 2: On bounded $M$:}\\
\textbf{Step 3: Unbounded $M$:}\\
\textbf{Step 1:}\\
Since $X$ a Hilbert space and $M$ a bounded convex closed set, we define
\begin{align*}
	\tilde{F}: M \rightarrow X^*,\quad u \mapsto \langle R^{-1} F(u), \cdot \rangle 
\end{align*}
where $R$ the Riesz map. By the approximation theorem in the Hilbert space we have for $w \in X$
\begin{align*}
	\langle w - P_M w,  v - P_M w \rangle \le 0 \quad \forall v \in M
\end{align*}
Note if we have $P_M w = u \in M$ and $w = u - \tilde{F}(u) $, we would be able to conclude the variational inequality. To prove that, observe the 2 equation above defines a fixed point map:
\begin{align*}
	G: M \rightarrow M, \quad G =  P_M(I - \tilde{F}) 
\end{align*} 
If $G$ admits a fixed point, then we have
\begin{align*}
	&G(u) = u\\
	&\Leftrightarrow
	u = P_M u - P_M\tilde{F}(u)\\
	& \Rightarrow P_M w = u, \quad w - P_M w = - \tilde{F}(u)
\end{align*}
Now we prove the fixed point exists by Brouwer, we need to prove $F$ therefore $G$ is continuous: By the boundedness of $M$ $u_n$ is bounded, by \textbf{(F1)} $F(u_n)$ is bounded in $X$.\\
$\Rightarrow$: $F(u_n)$ admits a convergent subsequence $u_k$ since $\dim X < \infty$, $F(u_k) \rightarrow u^*$. $u^*$ is a limit therefore also a weak lower limit. \\
$\Rightarrow$: $\langle F(u) - u^*, u - v \rangle \le 0 \quad \forall v \in B_\epsilon M$ by \textbf{(F3)}. \\
$\Rightarrow$: Choose $v = u \pm \epsilon v$ we have $F(u) = u^*$
$\Rightarrow$: $F$ is continuous hence $G$. \\
\textbf{Step 2:} Define $M_\delta := M \backslash B_\delta(\partial M)$, since $\dim X < \infty$, by the boundedness of $M$, there exists a convergent subsequence for $u_{\delta_k}$ in $X$, $F(u_{\delta_k})$ in $X^*$:
\begin{align*}
	\texttt{as }\delta_k \downarrow 0, \quad u_{\delta_k} \rightarrow u \texttt{ in } X, \quad
	F(u_{\delta_k}) \rightarrow u^* \texttt{ in } X^* 
\end{align*}
Furthermore, by step 1 we have
\begin{align*}
	\langle F(u_{\delta_k}), u_{\delta_k} - v \rangle \le 0 \quad \forall v \in M
\end{align*}
hence
\begin{align*}
	\langle u^*, u - v \rangle \leftarrow \langle F(u_{\delta_k}), u_{\delta_k} - v \rangle \le 0
\end{align*}
By the assumption \textbf{F3} we have
\begin{align*}
	\langle F(u), u - v\rangle \le \langle u^*, u - v \rangle \le 0
\end{align*}
\textbf{Step 3: }
$\forall R \in \mathbb{N}$, define $M_R:= \{ v \in X: \|v\|_X \le R\}$, by step 2 we have
\begin{align*}
	\forall R  \exists u_R \in M_R: \langle F(u_R), u_R - v \rangle \le 0 \qquad \forall v \in M_R
\end{align*}
and by the coercivity we know for all $u_0$ there exists some $C_0$, such that
\begin{align*}
	\langle F(u), u - u_0 \rangle > 0  \quad \forall \|u \| \ge C_0
\end{align*}
choose any large $R \ge \max\{C_0, \|u_0\|\}$, in particular $u_0 \in M_R$, we know by the coercivity the $u_R$ we get by step 2 must satisfies
\begin{align*}
	\|u_R\| \le C_0
\end{align*}
$\forall v \in M$: by Minty's trick contracting $v \in M$ and make  $v_\epsilon = (1 - \epsilon) u + \epsilon v \in M \cap B_\epsilon(u_R) $, in particular $v_\epsilon \in M_R$ for $\epsilon$ small enough, we have
\begin{align*}
	0 \ge \langle F(u_R), u_R - v_\epsilon \rangle = \epsilon \langle F(u_R), u_R - v \rangle \quad \forall v \in M
\end{align*} 
\end{proof}

\begin{theorem}
	Given $X$ a Hilbert space and $F: M \rightarrow X^*$ satisfies \textbf{F1-F3}. The variational inequality admits a solution. 
\end{theorem}
\begin{proof}
	By Galerkin approximation. The setup: a nested sequence of subspace $X_n \subset X_{n + 1} \subset X$. $u_0 \in X_1$, $\bigcup_n X_n = X, M_n = M \cap X_n$. The construction is possible since $X$ is separable. \\
	Define:
	\begin{align*}
		F_n: M_n \rightarrow X_n^*, \quad \langle F_n(u), v \rangle = \langle F(u), v \rangle \qquad \forall v \in X_n
	\end{align*}
	\textbf{1):}Verify that $F_n$ satisfies the assumption \textbf{F1-3}.
	\begin{enumerate}
		\item $\|F_n\|_{X^*} \le \|F\|_{X^*}$, hence $F_n$ must also map the bounded set to bounded set. 
		\item $u_0 \in X_1 \Rightarrow u_0 \in M_n$.
		\item Suppose $u_n \rightharpoonup u, F_n(u_n) \rightharpoonup u^*$, since $\dim X < \infty$ the weak convergence implies strong convergence, we have
		\begin{align*}
			\langle F_n(u_n), u_n \rangle \rightarrow \langle u^*, u \rangle  
		\end{align*}
		For $\langle F_n(u) - u^*, u - v \rangle \forall v \in M_n$: Since $u_n \in M_n$ weakly(and strong) convergent in $X_n$  therefore  
		\begin{align*}
			\langle F_n(u_m), u_m) \rangle = \langle F(u_m), u_m) \rangle \rightarrow \langle \tilde{u}^*, u\rangle
		\end{align*}
		Since $\langle F_n(u_m), v \rangle = \langle F(u_m), v \rangle \le \langle u^*, v \rangle$
	\end{enumerate}
	\textbf{2): The weak compactness}
	By the coercivity condition \textbf{F2} we have
	\begin{align*}
		\|u_n\| \le C_0 \quad \forall n \in \mathbb{N}
	\end{align*}
	By \textbf{F1} we know $F$ maps bounded set to bounded set, therefore
	\begin{align*}
		F(u_n) \le \infty \quad \textbf{ in } X
	\end{align*}  
	by reflexivity of $X$ we have (without relabeling)
	\begin{align*}
		F(u_n) \rightharpoonup u^*
	\end{align*}
	\textbf{3): Passage to the limit}
	we have
	\begin{align*}
		\langle F(u_n) , u_n - v \rangle \le 0  \qquad \forall v \in M_n
	\end{align*}
	It follows
	\begin{align*}
		&\langle F(u_n), u_n \rangle \le \langle F(u_n), v \rangle \rightarrow \langle u^*, v \rangle \texttt{ in } M_\infty\\
		(\textbf{F3})
		& \Rightarrow \langle u^*, u \rangle \le   \liminf_n \langle F(u_n), u_n \rangle \le \langle u^*, v \rangle\\
		&\Rightarrow \langle u^*, u - v \rangle \le 0
	\end{align*}
	Choose $v = u$ we have
	\begin{align*}
		\langle u^*, u \rangle \le  \liminf \langle F(u_n), u_n \le \langle u^*, u \rangle 
	\end{align*}
	We have
	\begin{align*}
		\liminf_n \langle F(u_n), u_n \rangle = \langle u^*, u \rangle
	\end{align*}
	therefore by \textbf{F3}
	\begin{align*}
		\langle F(u) - u^*, u - v \rangle \le 0 \quad \forall v \in M
	\end{align*}
	therefore
	\begin{align*}
		\langle F(u), u - v \rangle \le \langle u^*, u - v \rangle \le 0
	\end{align*}
\end{proof}

\begin{example}[Elastic Beam]
	Observe
	\begin{align*}
		I(v) := \int_I \frac{a_0}{2} |v''| + f \cdot \nu  dx 
	\end{align*}
	with 
	\begin{align*}
		M:= \{ v \in H^2(I) \cap H^1_0(I) | \quad v \ge g \quad a.e. \quad g(0), g(l) \le 0 \}
	\end{align*}
\end{example}


\subsection{The p-Laplacian}
Consider for $\Omega \subset \mathbb{R}^d$
\begin{align*}
	\inf_{v \in W^{1, p}_0}I(v):= \int_\Omega \frac{1}{p}|\nabla v |^p + \frac{\mu v }{2} + fv dx 
\end{align*}
we arrive at the ELE:
\begin{align*}
	-\div (|\nabla u |^p \nabla u | + \mu u = f &\texttt{ in } \Omega\\
	u = 0 &\texttt{ on } \partial\Omega
\end{align*}



\subsection{Perturbed monotone operators}
In the perturbed case we have the operator consisting of 2 parts, usually one described by the non-linearity and one by linearity(or continuity). We consider the non-linear part the principle term and the linear part the perturbed term, and extend our previous result regarding the non-linear elliptic variational inequality. 
\begin{theorem}
	Given $A: M \times M \rightarrow X^*$ with the general setting. (M closed, convex, X reflexive). Define $F(u):= A(u, u)$, if $A$ satisfies
	\begin{enumerate}
		\item $A$ is monotone in the second argument.
		\item $\forall v \in M$:  $u_n \rightharpoonup u \Rightarrow A(u_n, v)$ in M $\overset{*}{\rightharpoonup} A(u, v)$ in $X^*$.
		\item $\forall v \in M$: $u_n \rightharpoonup u$ $\Rightarrow \langle A(u, v), u \rangle \le \liminf_{n \rightarrow \infty} \langle A(u_n, v), u_n  \rangle$.   
	\end{enumerate}
\end{theorem}
\begin{proof}
\end{proof}


\subsection{Quasi-linear Elliptic systems}
Consider
\begin{align*}
	- \div (a(u, \nabla u) ) + f(u) = 0 & \texttt{ in } \Omega\\
	u = 0 & \texttt{ on } \partial \Omega 
\end{align*}

We assume:
\textbf{The Growth condition}
\begin{enumerate}
	\item $\forall \delta > 0 \exists C_\delta$:
	\begin{align*}
		a(z, q) \le C_\delta (|q|^{p -1} + 1) + \delta |z|^{p -1}
	\end{align*}
	\item There exists $C > 0$ s.t.
	\begin{align*}
		f(z) \le C |z|^{p - 1}
	\end{align*}
\end{enumerate}




\section{Classical Maximum Principle}
\subsection{Weak Maximum Principle}

\subsection{Harnack's Inequality}
\begin{theorem}
	Given $u$ a solution( i.e., $\partial_t u + Lu = 0$) with $u \ge 0$ in $\Omega_T$ and $u \in C^\infty(\Omega_T)$. Assuming $a_{ij} = a_{ji} \in C^\infty(\Omega_T), b_i = 0, c_i = 0$, $\Omega' \subset \subset \Omega$ subdomain,  then there exists $\forall t_1, t_2$ with  $0 < t_1 < t_2 \le T$ a constant $C$ such that
	\begin{align*}
		\sup_{\Omega'} u(t_1, \cdot) \le C \inf_{\Omega'} u(t_2, \cdot)
	\end{align*}
\end{theorem}

\subsection{Strong Maximum Principle}
\begin{theorem}[Strong maximum principle]
The coefficients of $L$ defined as before. 
Given $u$ a sub-solution(i.e., $\partial_t u + L u \le 0$) and $u \in C^\infty(\Omega_T) \cap C(\overline{\Omega_T})$. If $u$ admits its maximum at $(t_0, x_0) \in \Omega_T$, then 
\begin{align*}
	u \equiv const \texttt{ in } \Omega_{t_0}:= (0, t_0] \times \Omega
\end{align*}
\end{theorem}
\begin{proof}
Define $\Omega' \subset\subset$ with $(t_0, x_0) \in \Omega''$.
	Define $w$: 
	\begin{align*}
		&\partial_t w + Lw = -(\partial_t u + L u) \ge 0 \quad (t, x) \in \Omega_T'\\
		& w = 0 \quad \{0\} \times \Omega'
	\end{align*}
	The parabolic regularity asserts $w$ admits the same regularity. Then $v = w + u$ satisfies
	\begin{align*}
		& \partial_t v + Lv = 0 \quad (t, x) \in \Omega_T'\\
		& v = 0 \quad \{0\} \times \Omega'
	\end{align*}
	by the weak maximum principle we have
	\begin{align*}
		u \le v \le u \le u(t_0, x_0) := M
	\end{align*}
	Further define
	\begin{align*}
		\tilde{v}: = M - v
	\end{align*}
	we have
	\begin{align*}
		&\partial_t \tilde{v} + L\tilde{v} = 0 \quad \Omega' \\
		&\tilde{v} \ge 0 \quad \Omega'
	\end{align*}
	Apply Harnack's inequality we have
	\begin{align*}
		&\sup_{t \le t_0} \tilde{v} \le \tilde{v}(t_0, \cdot) = 0\\
		&\Rightarrow \tilde{v} \equiv 0 \quad \forall t \in (0, t_0)\\
		&\Rightarrow v \equiv const
	\end{align*}
	approximate $\Omega$ by $\Omega'$ we get the result. 
\end{proof}
 
 

\section{Reaction-Diffusion System}
\subsection{Global Existence of Lipschitz nonlinearity}
\begin{theorem}
Given the PDE
\begin{align*}
	\partial_t u - \Delta u = f(u) \quad &\texttt{ in } (0, T] \times \Omega\\
	u(0, ) = u_0 \quad &\texttt{ in } \{0\} \times \Omega\\
	u \equiv 0 \quad &\texttt{ 0n } (0, T] \times \partial\Omega
\end{align*}
where $f$ Lipschitz continuous with
\begin{align*}
	|f(u) - f(v)| \le C |u - v|
\end{align*}
Then there exists a unique solution 
\begin{align*}
	u \in C([0, T], L^2(\Omega))
\end{align*}
\end{theorem}
\begin{proof}
	proof with the contracting map. \\
	\textbf{Step 1:} Transform to a linear problem: define $h = f(u)$ and observe the system:
	\begin{align*}
		\partial_t w - \Delta w &= h\\
		w(0, ) &= u_0
	\end{align*}
	by the linear theorem there exists a \textbf{unique} solution 
	\begin{align*}
		w \in W(0, T; V, H) \hookrightarrow C([0, T], L^2)
	\end{align*}
	we can therefore define the map
	\begin{align*}
		\mathcal{A}: X_t \rightarrow X_t , \mathcal{A}(u) = w
	\end{align*}
	\textbf{Step2:} Show $\mathcal{A}$ is a contraction map
	observe the weak form:
	\begin{align*}
		\int_\Omega w_t \varphi + \mathbb{A} \nabla w \cdot \nabla \varphi dx = \int_\Omega f(u) \varphi dx \qquad \forall_{a.a.} t \in [0, T]
	\end{align*}
	Given $w, \tilde{w} \in \ran(\mathcal{A})$ and test with we have
	\begin{align*}
		\langle \partial_t(w - \tilde{w}), w - \tilde{w} \rangle 
		+ \|\nabla (w - \tilde{w})\|_{L^2} 
		&= \langle f(u) - f(\tilde{u}), w - \tilde{w} \rangle\\
		\frac{d}{dt}\|w - \tilde{w}\|_{L^2}^2 + C_p \|w - \tilde{w}\|_{L^2}^2 
		&\le C \|u - \tilde{u}\|_{L^\infty} \|w - \tilde{w}\|_V\\
	\end{align*}
	where the bracket indicates the linear mapping $\langle , \rangle_{V^*, V}$. By the Young's inequality the righthand side is bounded by 
	$C_\delta( \|u - \tilde{u}\|_{L^\infty}^2 + \delta \|w - \tilde{w}\|^2_{V}$, and again by Poincare the $V$ norm could be written as $L^2$ norm up to a constant. In summary we have after integration over $t$
	\begin{align*}
		\|(w - \tilde{w})(t)\|_H^2 \le t C \|u - \tilde{u}\|_{L^\infty}^2
	\end{align*}
\end{proof}





\subsection{The Heat kernel}



\subsection{Local well-posedness for locally Lipschitz nonlinearity}


\subsection{Maximum Life Span}


\subsection{Blow-up}



\section{Hyperbolic Equation}
We observe the problem:
\begin{equation}
	\partial_t u + \sum_{i = 1}^d A_j \partial_j u = B u + f
\end{equation}
where $u:[0, T) \times \mathbb{R}^d \rightarrow \mathbb{R}^n$ for $T \in (0, \infty]$ with $A_j(t, x), B(t, x) \in C_b^\infty$. We introduce the notation:
\begin{align*}
	&L = L(t) = \sum_{j = 1}^d A_j(t) \partial_j x - B(t)\\
	&Pu = \partial_t u + Lu
\end{align*}


\subsection{fractional Sobolev spaces}
Recall the definition of the Schwartz space:
\begin{definition}[Schwartz space]
\begin{equation}
	\mathcal{S}(\mathbb{R}^d):= \{ f \in C^\infty(\mathbb{R}^d, \mathbb{C}): |x^\alpha \partial^\beta u(x) \le C_{\alpha, \beta}, \texttt{for all } x \in \mathbb{R}^d\}	
\end{equation}
\end{definition}

In the Fourier transform we know the differentiation becomes multiplication of the monoid, i.e.


\subsection{Constant coefficient system}
Observe the problem
\begin{align*}
	& \partial_t u + \sum_{ j = 1}^d A_j \partial_j u = 0 \texttt{ in } (0, T) \times \mathbb{R}^d\\
	& u(0, \cdot) = u_0 \texttt{ in } \mathbb{R}^d
\end{align*}


\subsection{Linear symmetric hyperbolic system}
Observe the system
\begin{align*}
	& \partial_t u + \sum_{j = 1}^d A_j \partial_j u = Bu + f(t, x) \texttt{ in } (0, T)\\
	& u(0, \cdot) = u_0 \texttt{ in } \mathbb{R}^d
\end{align*}
The coefficients are assumed to be bounded and smooth with bounded derivatives. We now define 
\begin{align*}
	& L(t) = \sum_{j = 1}^d A_j \partial_j u - B\\
	& P(u) = Lu + \partial_t u
\end{align*}

\begin{definition}[The hyperbolic system]
\end{definition}

\begin{definition}[The symmetric hyperbolic system]
\end{definition}



\subsection{The energy estimate}
\begin{definition}[The distributional solution]
\end{definition}
why do we need the concept of distributional solution. 

\begin{proposition}[The energy estimate($L^2$ energy estimate)]
	Given $T > 0$ and $S_0(x) \in C_b^\infty(\mathbb{R}^d, \mathbb{R}^{n \times n})$. with $0 < \lambda \le \Lambda$ s.t. 
	\begin{align*}
		\lambda \le S_0(x) \le \Lambda \quad \forall x \in \mathbb{R}^d,
	\end{align*}
	then we have the  
\end{proposition}
and the equation may be rewritten as 
\begin{equation}
	Pu = f
\end{equation}

\begin{proposition}[The general energy estimate]
\end{proposition}







